\section{Propagation of oscillatory packets}\label{sec:packets}

The principal result of this section is the uniform estimate in
Theorem~\ref{pkt:main} for arbitrary finite complex packet arrays and
coordinate masks.  We define its exact operators and atomic norms,
then prove the composition, localization, and fourth-power estimates
used in Section~\ref{sec:extremal}.  That section derives the conditional
time profile \eqref{eq:source-1} required by geometric propagation from
near-extremal packet arrays themselves.

\subsection{The phase and an exact packet frame}

The Fourier-series construction and changes of packet frame follow the
established wave-packet framework; see
\cite[arXiv:math/0210084v2, Section~4 and Lemma~4.1]{Tao2003} and
\cite[arXiv:2104.11188v1, Section~4.1]{GOWWZ2025Packets}.
We prove the exact frame identities and coefficient estimates used here.

A finite partition of the sphere into sufficiently small graph patches
reduces the local analysis to the phase
\[
 x\cdot\xi+t\phi(\xi),\qquad \xi\in\R^2,
\]
with a smooth elliptic function $\phi$; the surface measure factor is
absorbed into the data.  We fix such a patch and phase throughout this
section.  We may extend $\phi$ to all of $\R^2$ so that its Hessian is
uniformly definite and bounded, and all its derivatives of order at
least two are bounded.  Indeed, on a patch of radius $\rho$, replace
$\phi$ outside the patch by its quadratic Taylor polynomial, using a
cutoff supported on the patch of radius $2\rho$.  The resulting change
in the Hessian is $O(\rho)$, since the Taylor remainder and its first two
derivatives are $O(\rho^3)$, $O(\rho^2)$, and $O(\rho)$, respectively.
Taking $\rho$ small preserves definiteness.  A reflection of the time
coordinate exchanges the positive and negative definite cases.  Write
\[
 U_\xi=-\nabla\phi(\xi).
\]
On each fixed bounded frequency region, $\xi\mapsto U_\xi$ is
quantitatively bi-Lipschitz.  All constants below may depend on the
fixed phase, the frequency region, and the chosen windows.

Once this continuation is fixed, the arrays below may use frequency
labels anywhere in a fixed bounded region, including the transition
region of the cutoff.  Their windows need not lie in the original
spherical patch.  The phase and all its derivative bounds remain fixed
before the packet scales vary.  The global sphere argument will return
to data supported on the original graph patches in
Section~\ref{sec:global}.

Let $N$ be dyadic, put $\delta=N^{-1}$, and use the oscillatory factor
\[
 e^{2\pi iN^2(x\cdot\xi+t\phi(\xi))},\qquad 0\le t<1.
\]
Choose a smooth nonnegative function $\chi$ whose integer translates
form a square partition of unity.  In particular, with
\[
 \chi_\theta^\nu(\xi)=\chi\bigl((\xi-\nu)/\theta\bigr),
 \qquad \nu\in\theta\Z^2,
\]
we have $\sum_\nu(\chi_\theta^\nu)^2=1$.  Fix $L_{\mathrm f}$ so that
each window is supported in a square of side $L_{\mathrm f}\theta$.
For a dyadic factor $1\le m\le N$, set
\begin{equation}\label{pkt:scales}
 \theta=m\delta,\qquad r=m^{-2},\qquad
 w=r\theta=\frac{\delta}{m}.
\end{equation}
Partition $[0,1)$ into intervals $I$ of length $r$, and let $t_I$ be the
left endpoint of $I$.  A packet index at this scale is
$v=(\nu_v,z_v,I_v)$, where
\[
 \nu_v\in\theta\Z^2,
 \qquad z_v\in L_{\mathrm f}^{-1}w\Z^2.
\]
At one time $t_I$, define the analysis and synthesis operators by
\begin{align}
 (\mathsf P_{m,I}f)_{\nu,z}
 &=\int f(\xi)\chi_\theta^\nu(\xi)
       e^{2\pi iN^2(z\cdot\xi+t_I\phi(\xi))}\,d\xi,
       \label{pkt:analysis}\\
 \mathsf S_{m,I}d(\xi)
 &=(L_{\mathrm f}\theta)^{-2}
    \sum_{\nu,z}d_{\nu,z}\chi_\theta^\nu(\xi)
       e^{-2\pi iN^2(z\cdot\xi+t_I\phi(\xi))}.
       \label{pkt:synthesis}
\end{align}
Synthesis is initially defined for finite arrays, and then by continuity
on $\ell^2$.

\Needspace{4\baselineskip}
\begin{lemma}[Frame identity]\label{pkt:frame}
For each $I$,
\[
 \mathsf S_{m,I}\mathsf P_{m,I}f=f,
 \qquad
 \|\mathsf P_{m,I}f\|_{\ell^2}^2
   =(L_{\mathrm f}\theta)^2\|f\|_2^2,
 \qquad
 \|\mathsf S_{m,I}\|_{\ell^2\to L^2}
   =(L_{\mathrm f}\theta)^{-1}.
\]
Consequently, if $I_+\subset I_-$ and
$\ell=m_+/m_-$, the transition
$\mathsf P_{m_+,I_+}\mathsf S_{m_-,I_-}$ has
$\ell^2\to\ell^2$ norm at most $\ell$.
\end{lemma}

\begin{proof}
The identity $N^2w\theta=1$ shows that, as $z$ varies over its grid,
$N^2z$ varies over $(L_{\mathrm f}\theta)^{-1}\Z^2$.  Thus the analysis
entries for one frequency window are the unnormalized Fourier
coefficients on a square of side $L_{\mathrm f}\theta$.  Parseval gives
\[
 \sum_z|(\mathsf P_{m,I}f)_{\nu,z}|^2
  =(L_{\mathrm f}\theta)^2\|f\chi_\theta^\nu\|_2^2.
\]
Sum over $\nu$.  Fourier inversion, with its factor
$(L_{\mathrm f}\theta)^{-2}$, gives
$\mathsf S_{m,I}\mathsf P_{m,I}f
=f\sum_\nu(\chi_\theta^\nu)^2=f$.
Finally, $\mathsf S_{m,I}$ is
$(L_{\mathrm f}\theta)^{-2}$ times the adjoint of
$\mathsf P_{m,I}$.  The operator norm statements follow.
\end{proof}

The initial scale is $m=1$, with the single time interval $[0,1)$.
A \emph{single-step operator} from factor $1$ to factor $m$ consists of
synthesis at factor $1$, followed by analysis in every terminal time
bin, followed by any coordinate mask.  More generally, successive
transitions use synthesis independently in each parent time bin and
analysis in its child bins.  Each finite array has a specified retained
coordinate set, and an index is called \emph{active} when it belongs to
that set, even if its coefficient is zero.  A coordinate mask projects
onto a subset of the retained set.  If \(M_I\) denotes the terminal
coordinate mask, the single-step map is exactly
\[
 (T_m d)_I=M_I\mathsf P_{m,I}\mathsf S_{1,[0,1)}d.
\]
The initial array is arbitrary and need not belong to the range of an
analysis operator.  Every map retains the full transition matrix
between its active index sets.  The localization graphs introduced
later count possible interactions; they do not delete individual
matrix entries or summands.

\subsection{Capacity and the atomic norm}

Recall the biased area
$W(E)=L^{1+\kappa}a^{1-\kappa}$ for an ellipse with radii $L\ge a>0$.
For an array in one time bin at scale $(\theta,r,w)$, define
\begin{equation}\label{eq:source-29}
 \begin{split}
 K(c)&=\theta^2
   \sup_{\substack{E:\text{both radii at least }\theta}}
   \frac{1}{W(E)}
   \sum_{\substack{U_{\nu_v}\in u+E\\ z_v\in z+rE}}|c_v|^2,
       \\
 G(c)^3&=\left(w^2\sum_v|c_v|^2\right)K(c)^{1/2}.
 \end{split}
\end{equation}
The supremum includes both independent translations $u,z\in\R^2$.
There is no upper restriction on the radii.  Define the lattice atomic
norm by
\[
 A(c)=\inf\left\{\sum_{j=1}^qG(a_j):
          |c|\le\sum_{j=1}^q|a_j|,\ q<\infty\right\}.
\]
The inequality in this definition is coordinatewise.  For arrays over
all bins of length $r$, put
\begin{equation}\label{pkt:combined-norm}
 \mathcal A(c)=\left(r\sum_I A(c_I)^3\right)^{1/3}.
\end{equation}

\Needspace{4\baselineskip}
\begin{lemma}[Atomic norm properties]\label{pkt:atomic}
On each finite index set, $A$ is a norm, is monotone under absolute
coordinatewise domination, and is contracted by coordinate masks.
The same assertions hold for $\mathcal A$.  Moreover,
\begin{equation}\label{eq:source-30}
 r w^2\sum_v|c_v|^3\le\mathcal A(c)^3.
\end{equation}
An operator estimate with input $G(d)$ and output a norm remains true
with input $A(d)$ and the same constant, provided the estimate holds
for every input on the specified support.
\end{lemma}

\begin{proof}
A translated disk of radius $\theta$ can include any chosen entry in
both tests in \eqref{eq:source-29}.  Since its biased area is
$\theta^2$, we have
\[
 K(c)\ge\|c\|_\infty^2,
 \qquad
 G(c)^3\ge w^2\|c\|_2^2\|c\|_\infty
          \ge w^2\sum_v|c_v|^3.
\]
The triangle inequality in $\ell^3$ transfers this lower bound to
every atomic decomposition and hence to $A$.  A coordinate atom has
cost $w^{2/3}|c_v|$, so $A$ is finite and positive away from zero.
Homogeneity, the triangle inequality, and absolute monotonicity follow
directly from its definition.  The weighted $\ell^3$ sum in
\eqref{pkt:combined-norm} proves the corresponding statements for
$\mathcal A$ and gives \eqref{eq:source-30}.

Atoms in a decomposition can first be restricted to $\supp d$, since
$G$ is monotone.  If $|d|\le\sum_j|a_j|$, choose, at each coordinate,
complex numbers $d_j$ with the phase of $d$, with
$\sum_jd_j=d$ and $|d_j|\le|a_j|$.  For a linear map $T$ and an output
norm $\mathcal B$, a bound $\mathcal B(Tb)\le C G(b)$ therefore gives
\[
 \mathcal B(Td)\le\sum_j\mathcal B(Td_j)
       \le C\sum_jG(a_j).
\]
Taking the infimum proves the last assertion.  In particular, support
restrictions imposed on the input are inherited by all atoms.
\end{proof}

\subsection{Uniform exponents and composition}

\begin{definition}[Admissible families]\label{pkt:families}
An admissible family has $m\to\infty$, with $m,N$ dyadic and
$m\le N\le m^B$ for a fixed finite $B$.  Its arrays have finite
supports, all their position labels have size at most $m^B$, and all
frequency labels lie in a fixed bounded set.  Enlarging $B$ is allowed;
it must remain fixed along each family.  Constants controlling the
phase and windows are fixed as well.  The same convention applies to
diagrams with a fixed finite number of scales.
\end{definition}

The full position lattices are infinite, but their retained portions
have polynomial size.  At a factor $q\le m$, the frequency and position
bounds permit at most
\[
 C(N/q)^2(m^BqN)^2=Cm^{2B}N^4
\]
indices per bin, and there are $q^2\le m^2$ bins.  The norms are also
comparable to the largest coordinate up to polynomial factors.  For
example, if one bin has $D\ge1$ retained coordinates and
$M=\max_v|c_v|$, the singleton test and the decomposition into coordinate
atoms give
\[
 w^{2/3}M\le A(c)\le D w^{2/3}M.
\]
Likewise, $M^2\le K(c)\le D M^2$, since $W(E)\ge\theta^2$.
These comparisons and the weighted bin norm show that a coordinate
error bounded by an arbitrarily large negative power of $m$ times the
largest input coordinate is negligible in all the norms used here.

Let $\gamma$ be the supremum, over admissible families and subsequences,
of the limiting upper exponents in $m$ of
\begin{equation}\label{eq:source-31}
 \frac{\mathcal A(T_m d)^3}{G(d)^3},\qquad d\ne0,
\end{equation}
where $T_m$ is a single-step operator, including any terminal mask.
By Lemma~\ref{pkt:atomic}, upper bounds have the same meaning with
$A(d)^3$ in the denominator.  An exponent bound is uniform on every
fixed complexity range, with any additional positive power loss:
otherwise a sequence violating that bound would itself be an
admissible family.  We use $m^{o(1)}$ in precisely this sense.  In
particular, it does not assert uniformity when a complexity parameter
escapes to infinity along a single family.

The packet estimate to be proved in Section~\ref{sec:extremal} is the
following.

\begin{theorem}[Uniform propagation of finite packet arrays]\label{pkt:main}
Fix \(0<\kappa<1/10\), the continued phase, a bounded frequency region,
and the window system above.  For every fixed admissible complexity
range and every \(\varepsilon>0\), there is a finite constant \(C\) such
that
\begin{equation}\label{pkt:uniform-main}
 \mathcal A(T_m d)^3
   \le C m^{10\kappa+\varepsilon}G(d)^3
\end{equation}
for every finite complex initial array \(d\) and every terminal
coordinate mask in that range.  The same estimate holds with \(A(d)^3\)
on the right.  The constant is uniform in \(m,N,d\), and the masks; it
may depend on \(\kappa,\varepsilon\), the complexity range, the fixed
phase and frequency region, and the windows.  Equivalently,
\begin{equation}\label{eq:source-32}
 \gamma\le10\kappa.
\end{equation}
\end{theorem}

We prove the theorem in Section~\ref{sec:extremal}.  The next lemmas
give the exact composition and analytic estimates used there.

\begin{lemma}[Composition]\label{pkt:composition}
Fix rational numbers $0=a_0<\cdots<a_J=1$ and dyadic factors
$m_0=1$, $m_J=m$, with $m_j=m^{a_j+o(1)}$.  Put
$\ell_j=m_{j+1}/m_j$.  If normalized transition $j$ has cubed norm
bound $C_j\ell_j^{\beta_j+\varepsilon}$, with $C_j$ independent of
$m$ and $\varepsilon>0$ arbitrary, the composed cubed norm bound is
\[
 \prod_{j=0}^{J-1}C_j\ell_j^{\beta_j+\varepsilon}
   =m^{\sum_{j=0}^{J-1}(a_{j+1}-a_j)\beta_j
          +\varepsilon+o(1)}.
\]
This remains true with arbitrary intermediate coordinate masks.
A better single-step bound restricted to particular input and output
supports can be used at any transition whose supports satisfy that
restriction, provided it is hereditary under restriction of the input
support.
\end{lemma}

\begin{proof}
Consider a parent bin $I'$ at factor $m'$, of length
$r'=(m')^{-2}$ and left endpoint $t_0$.  Make the change
\[
 t=t_0+r'\widetilde t,\qquad
 z=r'\widetilde z,\qquad N'=N/m'.
\]
The factor $e^{2\pi iN^2t_0\phi}$ is removed by rebasing the data.
The phase is again $N'^2(\widetilde x\cdot\xi+
\widetilde t\phi(\xi))$.  At a descendant factor $m''$, the normalized
factor is $\widetilde m=m''/m'$.  Its frequency width is unchanged,
whereas its spatial width and time length become
\[
 \widetilde w=w/r',\qquad \widetilde r=r/r'.
\]
The spatial grids consequently become exactly the grids prescribed
above.  The capacity is unchanged: the condition
$z_v\in z+rE$ becomes
$\widetilde z_v\in\widetilde z+\widetilde rE$.  Thus
\begin{equation}\label{pkt:normalization}
 \widetilde K=K,\qquad
 \widetilde G^3=(r')^{-2}G^3,\qquad
 \widetilde A^3=(r')^{-2}A^3.
\end{equation}
The last identity follows by applying the same common scalar factor
to every atomic cost.

An estimate with normalized cubed norm constant $C_{m''/m'}$ becomes
\[
 \frac{r''}{r'}\sum_{I''\subset I'}A(c_{I''})^3
       \le C_{m''/m'}A(d_{I'})^3.
\]
Multiply by $r'$ and sum over parent bins.  This gives the desired
global cubed bound at that transition.  Iterate it, using mask
contraction at every stage.  Taking logarithms in base $m$ gives the
stated weighted sum of exponents.  In particular, a fixed decrease in
$\beta_j$ decreases the full exponent by that amount times the positive
logarithmic length $a_{j+1}-a_j$.  At each such fixed split, all normalized
position and frequency ranges still have fixed polynomial complexity
in that step's scale ratio.  The uniformity in
Definition~\ref{pkt:families} therefore applies simultaneously to all
bins.  Finally, the last assertion follows from the support-preserving
atomic decomposition in Lemma~\ref{pkt:atomic}.
\end{proof}

\begin{lemma}[Interaction and finite truncation]\label{pkt:locality}
Let $v$ be a parent index and $b$ a child index, with $I_b\subset I_v$.
The transition matrix vanishes unless
$|\nu_b-\nu_v|\le C\theta_+$, and for every integer $L\ge0$ it obeys
\begin{equation}\label{pkt:matrix-decay}
 |T_{bv}|\le C_L
 \left(1+
  \frac{|z_b-z_v-(t_{I_b}-t_{I_v})U_{\nu_v}|}{w_-}
 \right)^{-L}.
\end{equation}
Consequently, for every fixed $\eta>0$, interactions outside the
enlarged parent beam of radius $m^\eta w_-$ give an error smaller than
any prescribed negative power of \(m\), times the largest input
coordinate, on an admissible family.  One may
also insert intermediate frame identities and truncate all intermediate
position grids to sufficiently large polynomial ranges, with the same
type of error.
\end{lemma}

\begin{proof}
The matrix entry is
\[
 (L_{\mathrm f}\theta_-)^{-2}
 \int\chi_{\theta_-}^{\nu_v}(\xi)
      \chi_{\theta_+}^{\nu_b}(\xi)
 e^{2\pi iN^2((z_b-z_v)\cdot\xi+
                    (t_{I_b}-t_{I_v})\phi(\xi))}\,d\xi.
\]
Disjoint frequency supports give the angular assertion.  On overlap,
rescale $\xi=\nu_v+\theta_-\zeta$.  The amplitude and all its
derivatives are uniformly bounded, since $\theta_-\le\theta_+$.
After the constant and linear phase terms are removed, the phase
derivatives are bounded because
\[
 N^2\theta_-^2(t_{I_b}-t_{I_v})
    =\frac{t_{I_b}-t_{I_v}}{r_-}\in[0,1).
\]
The remaining linear frequency is
$(z_b-z_v-(t_{I_b}-t_{I_v})U_{\nu_v})/w_-$.
Repeated integration by parts proves \eqref{pkt:matrix-decay}.

For finite retained arrays, sum the decay bound over the polynomial
number of entries counted after Definition~\ref{pkt:families}.  An
inserted frame identity instead uses the full position lattice; its
infinite tail is controlled by summing the same decay over lattice
annuli.  Fix the number of transitions, the admissible complexity, and
the positive beam-radius power first.  Then enlarge the intermediate
position ranges by fixed polynomial powers, and choose the
integration-by-parts order $L$ for the prescribed final error.
Frequency labels spread by only a bounded amount at each transition,
so a fixed number of insertions and truncations has negligible total
error.  The polynomial norm comparisons transfer this statement to
all the norms in use.  This estimates the error of a finite approximation
to an exact matrix product; it does not replace a transition by an
arbitrarily truncated set of edges.
\end{proof}

\Needspace{5\baselineskip}
\begin{lemma}[Preliminary bounds]\label{pkt:crude}
For a single-step map and each terminal bin,
\begin{equation}\label{pkt:mass-capacity}
 w^2\sum_{v\in I}|c_v|^2\le\delta^2\sum_v|d_v|^2,
 \qquad K(c_I)\le m^{4+o(1)}K(d).
\end{equation}
In particular, $\gamma\le2$.
\end{lemma}

\begin{proof}
The mass estimate is Lemma~\ref{pkt:frame}, since $w=\delta/m$.
For the capacity estimate, fix an output test ellipse $E$, with both
radii at least $\theta=m\delta$, and choose an arbitrarily small fixed
$\eta>0$.  By Lemma~\ref{pkt:locality}, the input entries contributing
to this test, apart from negligible error, have velocities in a
translate of $C m^\eta E$.  Their initial positions lie in another
translate of the same enlarged ellipse.  Indeed, a contributing
position equals its child position minus elapsed time times its
velocity, up to $m^\eta\delta$; the elapsed time is at most one and
both radii of $E$ are at least $\delta$.  Both translations are allowed
in the definition of $K(d)$.

Restrict the input to this one product of two translated tests.  Its
squared sum is at most
\[
 \delta^{-2}K(d)W(Cm^\eta E)
       \le C m^{2\eta}\delta^{-2}K(d)W(E).
\]
Apply the global squared operator norm bound $m^2$ to this restricted
input, and then multiply by the output factor $\theta^2/W(E)$.  Since
$\theta^2=m^2\delta^2$, the result is
$C m^{4+2\eta}K(d)$, with negligible error.  The estimate is uniform
over translations, orientations, and eccentricities.  Tail errors are
uniform too: $\theta^2/W(E)\le1$, and $K(d)\ge\|d\|_\infty^2$.
Let $\eta$ be arbitrarily small to obtain the stated exponent bound.

Finally, $A(c_I)\le G(c_I)$.  Multiply the mass bound by the uniform
upper bound for $K(c_I)^{1/2}$ and sum in time with weight $r$.
There are $r^{-1}$ bins, giving
$\mathcal A(c)^3\le m^{2+o(1)}G(d)^3$.
\end{proof}

\subsection{A fourth-power estimate from decoupling}

We use the ordinary \(\ell^2\) decoupling theorem of Bourgain and
Demeter~\cite[arXiv:1403.5335v3, Theorem~1.1, equation~(2)]{BourgainDemeter2015}.
For a smooth elliptic graph, Fourier thickness \(S^{-1}\), and graph
caps of diameter \(S^{-1/2}\), its critical three-dimensional estimate is
\[
 \Bigl\|\sum_\tau F_\tau\Bigr\|_4
 \le C_\epsilon S^\epsilon
          \Bigl(\sum_\tau\|F_\tau\|_4^2\Bigr)^{1/2}.
\]
The norms here are on all of spacetime. Uniformity over the normalized
phases below follows from the proof for general surfaces in
\cite[Section~7]{BourgainDemeter2015}: the ellipticity bounds are fixed,
and the common third-derivative bound controls the Taylor approximation
on each rescaled cap.

We derive the packet refinement for the precise class used here. The
induction follows Guth--Iosevich--Ou--Wang
\cite[arXiv:1808.09346v1, proof of Theorem~4.2]{GIOW2020}; compare the
weighted extension formulation in
\cite[arXiv:2411.08871v3, Section~4, Theorem~4.4]{WangWu2024}.
Guth--Iosevich--Ou--Wang also record an independent discovery of
the refinement by Du and Zhang. The proof below specifies which tubes count packet activity, including
at the smaller induction scale. The distinction between a packet's
activity region and its counted tube is substantive; a planar
counterexample to a support/counting mismatch is given in
\cite[arXiv:2608.28711v1, Proposition~2.1]{SinghSingh2026Refined}.

\begin{theorem}[Refined decoupling for fixed-profile packets]
\label{pkt:refined-decoupling}
Fix a nonzero, nonnegative \(a\in C_c^\infty(\R^2)\), and a smooth
elliptic phase \(\Phi\) on a fixed compact frequency region. Its
ellipticity and derivative bounds are fixed. A canonical packet at
scale \(R\ge2\) is
\[
 F_T(X,t)=d_T R\int
     a\bigl(\sqrt R(\xi-\xi_T)\bigr)
     e^{2\pi i((X-z_T)\cdot\xi+t\Phi(\xi))}\,d\xi .
\]
The frequency centers range over a fixed bounded set and lie on a grid
of spacing comparable to \(R^{-1/2}\). The compact phase region is
chosen to contain these centers and all their packet supports
\(\xi_T+R^{-1/2}\supp a\).
For each frequency center, the position centers lie on a grid of
spacing comparable to \(\sqrt R\). Assume the
position centers lie in a translate of a box of side \(C R\),
where \(C\) is fixed. Subsets of these grids are allowed.
The coefficients \(d_T\) are arbitrary complex numbers.

Fix \(0<\delta_0<1/8\). The counted tube associated to the packet is
\[
 T=\bigl\{(X,t): |t|\le3R,\quad
 |X-z_T+t\nabla\Phi(\xi_T)|\le R^{1/2+\delta_0}\bigr\}.
\]
The packet has transverse core width \(\sqrt R\), but need not vanish
outside this tube.  On the time slab below, the coefficient-relative
Schwartz bound proved below controls its decay away from the centerline.
The enlarged tube is the region used for counting; at each induction
scale we use the same definition with that scale's core width.
Let \(Y\) be a finite union of cubes from a fixed grid of side
\(\sqrt R\), contained in \(\{|t|\le2R\}\). Suppose each such cube
meets at most \(k\ge1\) counted tubes. For every \(\varepsilon>0\),
\begin{equation}\label{pkt:external-decoupling}
 \Bigl\|\sum_TF_T\Bigr\|_{L^4(Y)}^4
 \le C_{\varepsilon,\delta_0}R^\varepsilon
                         k R^2\sum_T|d_T|^4 .
\end{equation}
All constants may depend on the fixed profile, grid bounds, position
box constant, and phase bounds. The estimate is uniform over finite
subsets and the coefficients. If the weighted packet norms on a box
\(B\) satisfy
\(\|F_T\|_{L^4(w_B)}^4\asymp |d_T|^4R^2\), the right-hand side may
equivalently be written as
\(C_{\varepsilon,\delta_0}R^\varepsilon k
\sum_T\|F_T\|_{L^4(w_B)}^4\).
\end{theorem}

The proof groups the packets into coarser frequency caps and position
squares, then rescales each group to scale $\sqrt R$. Its main task is
to compare the tubes counted at that smaller scale with the original
multiplicity bound $k$. The local decoupling estimates can then be
summed with each packet charged through its fourth-power cost.

\begin{proof}
The grids give at most \(C R^2\) packets. Write
\(\mathcal C_R=R^2\sum_T|d_T|^4\).
For \(|t|\le3R\), integration by parts in the normalized cap gives,
for every fixed \(L\),
\[
 |F_T(X,t)|
 \le C_L|d_T|
 \left(1+\frac{|X-z_T+t\nabla\Phi(\xi_T)|}{\sqrt R}\right)^{-L}.
\]
Outside the counted transverse radius this gains \(R^{-L\delta_0}\).
The polynomial packet count can therefore be included before choosing
$L$ to obtain any prescribed power of decay.  Choose a fixed
band-limited Schwartz function \(\psi\), bounded away from zero on
\([-4,4]\), and put \(P_T=\psi(t/R)F_T\).
The spatial \(L^\infty\) bound is \(C|d_T|\), and Plancherel gives
\(\|F_T(\cdot,t)\|_2^2\le C|d_T|^2R\). Hence
\[
 \|P_T\|_4^4\asymp |d_T|^4R^2.
\]
For the lower bound, use a short interval \(|t|\le cR\) and a
ball of radius \(c\sqrt R\) about the core. The normalized integral
is bounded below there because \(a\) is fixed, nonnegative and
has positive integral. Thus no tail estimate is normalized by the
possibly small norm of an arbitrary oscillatory profile.

Let \(D(R)\) be the best constant in
\[
 \Bigl\|\sum_TF_T\Bigr\|_{L^4(Y)}^4
                       \le D(R)k\mathcal C_R
\]
over the fixed classes just specified. For the induction, strengthen
the allowed time slab slightly to
\(|t|\le(2+R^{-1/16})R\); keep the counted tubes in \(|t|\le3R\).
The same core bounds apply. This harmless buffer accommodates cells
meeting the endpoints after rescaling: the added length is
\(O(R^{3/4}+R^{1/2+\rho})\), and, at \(r=\sqrt R\),
\[
 R^{-1/16}+O(R^{-1/4}+R^{-1/2+\rho})<r^{-1/16}
\]
for sufficiently large \(R\), once \(\rho<1/16\).
For bounded \(R\), the packet count, the preceding cost bound and
H\"older's inequality give a finite bound for \(D(R)\). The zero
array is immediate, so assume \(\mathcal C_R>0\).

Partition the frequency centers into caps \(\tau\) of diameter
\(R^{-1/4}\). For each \(\tau\), partition the time-zero position
centers into squares of side \(R^{3/4}\). Assign each packet once to
its frequency cap and position square; denote the resulting disjoint
packet groups by \(\mathcal W_l\), and set
\(P_l=\sum_{T\in\mathcal W_l}P_T\).
The centerlines in a group remain in a fixed enlargement of the
corresponding coarse cylinder, of radius \(R^{3/4}\) and length
\(O(R)\). At any \(\sqrt R\)-cube, only boundedly many such cylinders
per coarse cap are nearby. Far cylinders have separation comparable
to \(R^{3/4}\) from their \(\sqrt R\) cores and contribute arbitrarily
small relative errors. No packet is multiplied by a coarse spatial
cutoff. The fine cap supports lie in fixed enlargements of their
assigned coarse caps. These enlargements have bounded overlap; a
fixed coloring separates them before ordinary decoupling, at a
constant cost.

Here is the exact normalization of a group. If \(c_\tau\) is its
frequency center and \(z_l\) its position-square center, put
\[
 \xi=c_\tau+R^{-1/4}\zeta,\qquad
 X'=R^{-1/4}\bigl(X+t\nabla\Phi(c_\tau)-z_l\bigr),\qquad
 t'=R^{-1/2}t,\qquad r=\sqrt R.
\]
After removing a common unimodular factor and absorbing each
packet's constant phase into its coefficient, the packets have the
same form at scale \(r\), with the same profile \(a\) and
\(|d'_T|=|d_T|\). Indeed
\(R\,d\xi=r\,d\zeta\), their frequency and position meshes become
\(r^{-1/2}\) and \(\sqrt r\), and the new position centers have size
\(O(r)\). The new phase has derivatives
\(R^{(2-j)/4}D^j\Phi\) for \(j\ge2\), so it stays in the fixed
elliptic class. The cutoff becomes \(\psi(t'/r)\).
The physical Jacobian is \(R\); consequently the inner cost
\(|d_T|^4r^2\), after returning to the original coordinates, is
exactly \(|d_T|^4R^2\).

Apply the inductive estimate at scale \(r\). Its counted radius
\(r^{1/2+\delta_0}\) pulls back to the narrower radius
\(R^{1/2+\delta_0/2}\). A pulled-back inner cube has transverse
size \(R^{1/2}\) and longitudinal size \(R^{3/4}\). Directions in
one coarse cap differ by \(O(R^{-1/4})\), so their relative
displacement across that cube is \(O(\sqrt R)\).
Choose
\[
 0<\rho<\min(\delta_0/8,\,1/16).
\]
If a pulled-back inner cube \(P\) meets a translate of an outer cube
\(Q\) by distance at most \(C R^{1/2+\rho}\), then every inner counted
tube meeting \(P\) belongs to an original tube meeting \(Q\). The
transverse distance is bounded by
\[
 C\bigl(R^{1/2+\delta_0/2}+R^{1/2+\rho}+\sqrt R\bigr)
                                  <R^{1/2+\delta_0}.
\]
This is where the distinction between the core and the counted tube
is used. Each induction step rescales the core and chooses its own
counted radius; it does not rescale the parent's already enlarged tube.

For clarity, we give the local decoupling bookkeeping, including the
weights. For every outer cube \(Q\), choose a translate and dilate
\(\eta_Q\) of one band-limited Schwartz function which is bounded
below on \(Q\). Its spatial and temporal scale is \(\sqrt R\).
Applying ordinary global decoupling to
\(\eta_Q\sum_lP_l\) is legitimate: multiplication broadens Fourier
support by \(O(R^{-1/2})\), exactly the thickness at the coarse
decoupling scale \(S=\sqrt R\). Keep the boundedly many nearby
coarse cylinders per cap. Their fourth-power estimate is
\[
 \Bigl\|\sum_TF_T\Bigr\|_{L^4(Q)}^4
 \le C_\beta R^\beta
       \left(\sum_l\|\eta_QP_l\|_4^2\right)^2
       +\text{relative tail errors},
\]
where \(\beta>0\) is arbitrarily small.

Truncate each local integral to cubes \(Q+j\sqrt R\) with
\(j\in\mathcal J\subset\Z^3\), \(|j|\le C R^\rho\).
The set of shifts is the same for every \(Q\), and
\(\#\mathcal J\le C R^{3\rho}\). On the retained region
\(|\eta_Q|\le C\). For each group \(l\), keep only inner cubes meeting one of these
retained translates, with \(Q\) ranging over the cubes of \(Y\).
The time-buffer inequality above places every such whole inner cube
in the strengthened child slab. Partition these cubes dyadically
by the number of inner counted tubes they meet. Write
\(Y_{l,h}\) for the pulled-back union with between \(h\) and \(2h\)
tubes, where \(h\) ranges over \(O(\log R)\) positive dyadic values.
The zero-count cells give only relative tail errors by the core
bound. Set
\[
 A_{l,j,h}(Q)
   =\int_{Q+j\sqrt R}
             \mathbf1_{Y_{l,h}}\,|P_l|^4.
\]
If this is nonzero, the preceding incidence comparison supplies at
least \(h\) packets from \(\mathcal W_l\) whose outer tubes meet
\(Q\). The groups are disjoint, so there are at most \(Ck/h\)
such groups \(l\). Therefore
\[
 \left(\sum_l A_{l,j,h}(Q)^{1/2}\right)^2
                       \le \frac{Ck}{h}\sum_l A_{l,j,h}(Q).
\]
Two finite Cauchy--Schwarz sums over \(j,h\) cost at most
\(C R^{6\rho}(\log R)^2\).
For each fixed shift \(j\), the cubes \(Q+j\sqrt R\) are disjoint.
Sum their integrals before estimating them by an entire inner-cell
union:
\[
 \sum_Q A_{l,j,h}(Q)\le\int_{Y_{l,h}}|P_l|^4
 \le C D(\sqrt R)\,h R^2
                         \sum_{T\in\mathcal W_l}|d_T|^4.
\]
The last inequality is the child estimate and the exact Jacobian
calculation above; the bounded common cutoff costs only a constant.
The factor \(h\) cancels the preceding \(h^{-1}\), and summation
over the disjoint groups preserves \(\mathcal C_R\). This order
avoids counting one long inner cell once for each outer cube it meets.

All discarded terms are homogeneous in the coefficients. For example,
\(\#\{T\}\le CR^2\) and
\(|\sum_{T\in\mathcal W_l}P_T|^4
 \le(\#\mathcal W_l)^3\sum_{T\in\mathcal W_l}|P_T|^4\).
The Schwartz bound on \(\eta_Q\), summed over the cube grid outside
\(R^\rho Q\), then bounds the off-neighborhood error by
\(C_LR^{8-(4L-3)\rho}\mathcal C_R\), before harmless fixed powers.
Choose its decay order after \(\rho,\delta_0\) and the desired error.
The core bound treats far coarse cylinders and zero-count child cells
in the same way. Summing over space uses integrable decay; no bound
on the number of distant selected cubes is needed.
Thus, for every prescribed \(A>0\), these terms are at most
\(C_A R^{-A}\mathcal C_R\) at the current induction scale.

We have proved the recurrence
\[
 D(R)\le
 C R^{\beta+6\rho}(\log R)^2D(\sqrt R)+C_A R^{-A}.
\]
Choose \(\rho\) smaller if necessary, after \(\delta_0\) and
\(\varepsilon\), and then \(\beta\), so that
\(\beta+6\rho<\varepsilon/4\). At a sufficiently large fixed base,
the logarithm and the fixed constant fit within another
\(R^{\varepsilon/8}\). The resulting exponent, including the
child bound \(R^{\varepsilon/2}\), is less than
\(7\varepsilon/8\). Induction from that base gives
\(D(R)\le C_{\varepsilon,\delta_0}R^\varepsilon\). The additive errors are absorbed
at their own scale; no error at the last small scale is asserted
to be a rapid-decay error in the original scale. The assumption
\(k\ge1\) allows their absorption into \(k\mathcal C_R\).
This proves \eqref{pkt:external-decoupling}.
The weighted version is the stated packet-cost comparison.
\end{proof}

\begin{proposition}[Discrete fourth-power propagation]
\label{pkt:quartic}
Consider one consecutive transition in an admissible family, with
scale ratio $\ell=m_+/m_-=m^{\alpha+o(1)}$, where $\alpha>0$ is fixed.
Fix finite retained parent and child coordinate sets and an adjacency
power $\eta>0$.  Let $d$ be any complex array on the parent sets, and
let $M_{I_+}$ project onto the retained set in child bin $I_+$.  The
output in that bin is
\[
 c_{I_+}=M_{I_+}\mathsf P_{m_+,I_+}
                    \mathsf S_{m_-,I_-}d_{I_-},
 \qquad I_+\subset I_-.
\]
Suppose each active child index has at most $k\ge1$ active parents in
its parent time bin satisfying
\[
 |\nu_b-\nu_v|\le m^\eta\theta_+,
 \qquad
 |z_b-z_v-(t_{I_b}-t_{I_v})U_{\nu_v}|
       \le m^\eta w_-.
\]
Thus $k$ counts nearby pairs of retained coordinates, independently of
whether their coefficients vanish.  The displayed operator still uses
every matrix entry between the retained sets.  Then
\begin{equation}\label{eq:source-33}
 r_+w_+^2\sum_b|c_b|^4
       \le m^{o(1)} k\,r_-w_-^2\sum_v|d_v|^4.
\end{equation}
For each fixed complexity range, $\alpha$, and $\eta$, and every
$\varepsilon>0$, the factor $m^{o(1)}$ may be replaced by
$C_\varepsilon m^\varepsilon$ for all sufficiently large $m$ in the
window
\[
 m^{\alpha/2}\le\ell\le m^{3\alpha/2}.
\]
The constant and the lower threshold within this window are uniform.
A family's entry into the window may depend on the rate at which
$\log_m\ell$ tends to $\alpha$.
\end{proposition}

\begin{proof}
Fix a parent time bin and one child frequency window.  Only parent
windows overlapping this child window contribute.  Select these
parents, but keep their synthesis packets intact.  In particular, we
do not multiply individual parent windows by the child cutoff before
applying refined decoupling.

Let $\nu_+$ be the child frequency center and $t_-=t_{I_-}$.
Set
\[
 \xi=\nu_++\theta_+\zeta,\qquad
 T=\frac{t-t_-}{r_+},\qquad
 X=\frac{x-(t-t_-)U_{\nu_+}}{w_+}.
\]
The identities $N^2w_+\theta_+=N^2r_+\theta_+^2=1$ give the
normalized phase, after removing the constant and linear terms,
\[
 \Phi_+(\zeta)=\theta_+^{-2}
 \bigl(\phi(\nu_++\theta_+\zeta)-\phi(\nu_+)
                 -\theta_+\nabla\phi(\nu_+)\cdot\zeta\bigr).
\]
It has uniformly definite Hessian and bounded higher derivatives on
a fixed compact set.  In these units the parent slab has length
\[
 R=\frac{r_-}{r_+}=\ell^2,
\]
the parent frequency caps have diameter $O(\ell^{-1})$, and the parent
packets have transverse width $\ell$. More precisely, the
normalized synthesis packet has the canonical form in
Theorem~\ref{pkt:refined-decoupling} with the fixed profile
$a=L_{\mathrm f}^{-2}\chi$ and centers
\[
 \zeta_v=\frac{\nu_v-\nu_+}{\theta_+},\qquad
 Z_v=\frac{z_v}{w_+}.
\]
Here the constant phase $e^{-2\pi iN^2z_v\cdot\nu_+}$ is absorbed
into $d_v$, preserving its absolute value.  The position mesh is
$L_{\mathrm f}^{-1}\sqrt R$. The normalized velocities are
bounded, because all selected parent windows overlap the child window.
The synthesis factor in \eqref{pkt:synthesis} cancels the area of a
parent frequency window, so each parent packet has amplitude
$O(|d_v|)$, with no extra scale factor.

For completeness, localize time by a fixed Schwartz function
$\psi(T/R)$ whose Fourier transform is compactly supported and which
is bounded away from zero on an interval containing $[0,1]$ in its
argument.  Denote the resulting parent packets by $P_v$.  Their Fourier
support lies in an $O(R^{-1})$ neighborhood of the normalized graph,
over caps of diameter $O(R^{-1/2})$.  On $|T|=O(R)$, integration by
parts on the parent cap gives transverse Schwartz decay on scale
$R^{1/2}$.  For larger $|T|$, the decay width and constants grow only
polynomially in $1+|T|/R$, which is harmless against $\psi(T/R)$.
Consequently,
\begin{equation}\label{pkt:packet-fourth-cost}
 \|P_v\|_4^4\le C|d_v|^4R^2
                       =C|d_v|^4\ell^4.
\end{equation}
We next pass from the discrete child samples to a spacetime integral.
Write $F=\sum_vP_v$.  In the normalized variables, the child cutoff is
a fixed smooth frequency multiplier.  All spatial and temporal
frequencies of $F$ lie in a fixed compact set.  There is therefore a
Schwartz kernel $H$ on spacetime, with uniform bounds, whose Fourier
transform implements the child multiplier on that set.  An active
child index $b$ has normalized sample location
\[
 p_b=(X_b,T_b)=\left(
   \frac{z_b-(t_{I_b}-t_-)U_{\nu_+}}{w_+},
   \frac{t_{I_b}-t_-}{r_+}\right).
\]
At these retained coordinates, up to a unimodular factor,
\[
 \psi(T_b/R)c_b=(H*F)(p_b).
\]
The time coordinates are integers from $0$ to $R-1$; at each time,
the spatial coordinates lie in a translate of
$L_{\mathrm f}^{-1}\Z^2$.  Thus the retained samples have bounded
multiplicity at unit resolution.  H\"older's inequality and summation
over these samples give
\[
 \sum_b|c_b|^4
       \le C\int |F(q)|^4
                    \sum_b|H(p_b-q)|\,dq.
\]
The summed kernel is uniformly bounded and decays rapidly with the
distance to the active sample set.  Truncate this distance at $m^\beta$,
where $\beta>0$ will be chosen small.  The omitted integral is bounded by
an arbitrarily large negative power times
$R^2\sum_v|d_v|^4$: use the polynomial number of input packets,
\eqref{pkt:packet-fourth-cost}, and the triangle inequality.  Cover the
retained neighborhood by cubes of side $\ell$ and denote their union
by $Y$.  We have reduced the desired estimate to a bound for
$\int_Y|F|^4$.

We check carefully that the degree hypothesis controls the tube
multiplicity on these cubes.  Choose
$\beta<\min(\eta,\alpha)/10$.  Every cube in $Y$ has an associated active
sample within $O(\ell+m^\beta)$ of all its points.  An enlarged parent
tube meeting that cube has, at the time of this sample, transverse
distance at most
\begin{equation}\label{pkt:cube-adjacency}
 C\bigl(\ell R^{\delta_0}+\ell+m^\beta\bigr)
\end{equation}
from the sample: the normalized tube velocities are bounded.
Choose $0<\delta_0<1/8$ with $3\alpha\delta_0<\eta/10$.
In the stated ratio window, $R^{\delta_0}\le m^{3\alpha\delta_0}
\le m^{\eta/10}$.  Then
\eqref{pkt:cube-adjacency} is less than $m^\eta\ell$ for sufficiently
large $m$, which is exactly the allowed parent-beam radius in child
units.  Angular overlap gives a fixed multiple of $\theta_+$, also
less than the permitted $m^\eta\theta_+$.  Each such tube is therefore
one of the at most $k$ parents counted at the associated sample.
This uses the stated adjacency radius, without replacing it by a
larger power.

The sample times satisfy $0\le T_b<R$, so the cube union $Y$ lies in
\[
 -C(\ell+m^\beta)\le T\le R+C(\ell+m^\beta).
\]
In the stated ratio window,
$(\ell+m^\beta)/R\le m^{-\alpha/2}+m^{\beta-\alpha}\to0$.
Thus $Y$ is contained in $\{|T|\le2R\}$ for sufficiently large $m$,
uniformly in that window, as required by
Theorem~\ref{pkt:refined-decoupling}.

Partition space into boxes of side $R$, aligned with the cube grid,
over this time slab.  In each box retain the packets whose centerlines
meet a fixed enlargement of
it.  Their time-zero centers lie in a box of side $O(R)$ because the
normalized velocities are bounded.  The other packets have rapidly
decaying tails there.  Admissibility gives a polynomial number of
boxes meeting $Y$, so choose the integration-by-parts order after this
fixed complexity to make the total discarded contribution negligible.
Each relevant tube is charged in only boundedly many boxes, since its
length and displacement in the slab are $O(R)$.

Write
$P_v=\psi(T/R)\widetilde P_v$, where $\widetilde P_v$ is the
underlying canonical extension packet.  On the selected cubes,
\[
 \left|\sum_vP_v\right|^4
   \le \|\psi\|_\infty^4
                  \left|\sum_v\widetilde P_v\right|^4.
\]
Apply the coefficient-cost estimate of
Theorem~\ref{pkt:refined-decoupling} to these canonical packets in
each box.  Their position centers lie in a box of side $O(R)$, their
coefficients remain arbitrary complex numbers, and their counted-tube
multiplicity is at most $k$.  The common cutoff and bounded box
multiplicity therefore give
\[
 \sum_{\substack{b:\text{fixed child angle}\\
                      I_b\subset I_-}}|c_b|^4
     \le m^{o(1)} k\,\ell^4
          \sum_{\substack{v\in I_-:\text{window overlaps}\ \nu_+}}
                         |d_v|^4.
\]
All tail errors can be absorbed into this bound, since $k\ge1$.
Since $R\le m^{3\alpha}$ in the ratio window, the decoupling loss
can be made smaller than any prescribed positive power of $m$
uniformly there.  For a prescribed final error power, choose first the
error allowed by Theorem~\ref{pkt:refined-decoupling}, then its
localization power $\delta_0$ and the sampling power $\beta$ small enough to
satisfy the preceding constraints.
Their constants are fixed before taking $m\to\infty$.

Each parent window overlaps only boundedly many child windows.  Sum
over child angles and parent bins and multiply by $r_+w_+^2$.  The
scale identity
\[
 r_+w_+^2\ell^4=r_-w_-^2
\]
proves \eqref{eq:source-33}.
\end{proof}

\begin{remark}[Weighted packet costs]
In the normalization of the preceding proof, the canonical packets
also have the weighted norms appearing in comparisons with refined
decoupling formulations.  Write
$P_v=\psi(T/R)\widetilde P_v$ as above.  For a spatial box $B$ used
there, with center $X_B$, put $q_B=(X_B,0)$ and
\[
 w_B(q)\asymp(1+|q-q_B|/R)^{-K_0},\qquad q=(X,T),\quad K_0>4.
\]
For every packet retained in this box,
\[
 \|\widetilde P_v\|_{L^4(w_B)}^4
       \asymp |d_v|^4R^2
       \asymp \|P_v\|_4^4.
\]
Indeed, the frequency-integral bound and Plancherel give
$\|\widetilde P_v(\cdot,T)\|_\infty\lesssim|d_v|$ and
$\|\widetilde P_v(\cdot,T)\|_2^2\lesssim|d_v|^2R$.
Integration with the temporal decay of $w_B$ or $\psi(T/R)$ proves
the upper bounds.  For sufficiently small fixed $c>0$, the
nonnegative profile gives $|\widetilde P_v(X,T)|\gtrsim|d_v|$ on
\[
 0\le T\le cR,\qquad
 |X-Z_v+T\nabla\Phi_+(\zeta_v)|\le c\sqrt R.
\]
This is the short-core argument in the proof of
Theorem~\ref{pkt:refined-decoupling}.  The region has volume
comparable to $R^2$.  Both $|\psi(T/R)|$ and $w_B$ are bounded below
there: the time-zero centers of the retained packets lie within
distance $O(R)$ of $B$, even if a tube meets the box only at a later
time.  This proves the lower bounds.
\end{remark}

\begin{remark}[Vanishing adjacency powers]\label{pkt:vanishing-powers}
Later there will be an outer index $i$, a fixed number $J$ of
transitions, and positive adjacency powers $\eta_i\to0$.  For each
fixed $i$ the scale ratios satisfy
$\ell_j=m^{1/J+o(1)}$, so $R=\ell_j^2=m^{2/J+o(1)}$.
Choose the external error power tending to zero with $i$, and its
localization power $\delta_i$ still smaller, with
$2\delta_i/J<\eta_i/10$.  Choose the sampling power less than
$\min(\eta_i,1/J)/10$.  The cube comparison above then uses precisely
the radius $m^{\eta_i}w_-$, while the coloring and decoupling losses
are $m^{o_i(1)}$.  Here $o_i(1)$ denotes an exponent whose limiting
upper absolute value as $m\to\infty$, with $i$ and $J$ fixed, tends
to zero as $i\to\infty$.  We take the limits in that order.  The
phase, profile, complexity, and positive radius powers remain fixed
in each inner limit.  Constants may depend on $i$; since $J$ is fixed,
the same convention applies to the product of all transition losses.
\end{remark}
